\documentclass{article}
\usepackage[margin=1.25in]{geometry}
\usepackage{parskip}
\usepackage{fancyhdr}
\usepackage{enumitem}

\newcommand{\coursename}{MA 214: Applied Statistics}
\newcommand{\termname}{Fall 2026}
\newcommand{\activitydate}{Week 3}
\newcommand{\projecttitle}{Project 1 Launch and Outline}
\newcommand{\submissionplatform}{Gradescope or Blackboard}
\newcommand{\duedate}{Sunday, September 20 at 10:00 PM}

\pagestyle{fancy}
\fancyhf{}
\lhead{\coursename{} (\termname{})}
\rhead{\activitydate{}}
\cfoot{\thepage}
\renewcommand{\headrulewidth}{.4mm}
\newcommand{\blankline}{\underline{\phantom{XXXXXXXXXXXXXXXXXXXXXXXX}}}
\newcommand{\blankid}{\underline{\phantom{XXXXXXXXXXXXXXX}}}

\begin{document}
\noindent\textbf{Group:}\blankline{}\hfill\textbf{Section:}\blankid{}

\medskip
\noindent\textbf{Member 1:}\blankline{}\hfill\textbf{BUID:}\blankid{}

\medskip
\noindent\textbf{Member 2:}\blankline{}\hfill\textbf{BUID:}\blankid{}

\medskip
\noindent\textbf{Member 3:}\blankline{}\hfill\textbf{BUID:}\blankid{}

\medskip
\noindent\textbf{Member 4:}\blankline{}\hfill\textbf{BUID:}\blankid{}

\begin{center}
  {\Large \textbf{\projecttitle{}}}
\end{center}

\subsection*{Purpose}
Today your group will launch Project 1 and work directly on the \textbf{Deliverable 1 outline}. By the end of lab, your shared Google Doc should contain a complete first draft that your group can revise and submit.

\subsection*{Shared Google Doc Setup}
\begin{enumerate}[itemsep=.25em]
  \item Create one Google Doc for the group and name it \texttt{P1_GroupNumber_Outline}.
  \item Share the document with all four group members and the instructor/TA with editing access.
  \item Add all four names and BUIDs at the top. Use headings that match the Deliverable 1 outline.
  \item Work in the same document throughout lab; do not create four separate drafts.
\end{enumerate}

\subsection*{Group-of-Four Roles}
\begin{itemize}[itemsep=.25em]
  \item \textbf{Member 1 -- Project lead:} title, background, and motivation.
  \item \textbf{Member 2 -- Data lead:} data source, key variables, observations, and limitations.
  \item \textbf{Member 3 -- Analysis lead:} research questions, objectives, hypotheses, and initial analysis plan.
  \item \textbf{Member 4 -- Editor and scheduler:} timeline, task assignments, formatting, and final quality check.
\end{itemize}

\subsection*{In-Class Work}
\begin{enumerate}[label={\bfseries (\Alph*)},itemsep=.5em]
  \item \textbf{Project orientation.} Review the Project 1 overview and rubric together; confirm the required sections and final submission format.
  \item \textbf{Choose the data.} Agree on one data set that can support a statistical analysis and model application. Record its source and access information.
  \item \textbf{Draft the questions.} Write at least two possible research questions or objectives, then select the primary question.
  \item \textbf{Build the analysis plan.} Identify a response variable, candidate predictors, first summaries or plots, and any likely limitations.
  \item \textbf{Peer edit.} Each member reads the complete draft and adds one specific improvement or clarification.
  \item \textbf{Check-in.} Ask the instructor/TA to review the primary question, data source, and feasibility of the proposed analysis.
\end{enumerate}

\subsection*{Deliverable}
Submit the completed Project 1 outline through \submissionplatform{} by \duedate{}. Your shared Google Doc must contain:
\begin{itemize}[itemsep=.25em]
  \item project title and all four group members;
  \item background and motivation;
  \item research question(s), objectives, or hypotheses;
  \item data source, key variables, observations, and limitations; and
  \item timeline and roles for the remaining project work.
\end{itemize}
Before submitting, confirm that every member has contributed, the document is readable without comments, and the final link is accessible to the instructor/TA.

\end{document}
